\documentclass[a4paper]{article}

\usepackage[english]{babel}
%\usepackage[utf8x]{inputenc}
\usepackage{amsmath}
\usepackage{graphicx}
\usepackage[colorinlistoftodos]{todonotes}
\usepackage[all]{xy}
\usepackage{dsfont}
\usepackage{amssymb}
\usepackage{bbold}
\usepackage{amsthm}

\usepackage{tikz}
\usepackage{verbatim}
\usetikzlibrary{matrix}
\usetikzlibrary{cd}

\usepackage[autostyle=true]{csquotes}
\usepackage{pdfpages}
\usepackage{hyperref}
\usepackage{fancyhdr}
\usepackage{multicol}



\newtheorem{theorem}{Theorem}[section]
\newtheorem{proposition}{Proposition}[theorem]
\newtheorem{lemma}{Lemma}[theorem]
\newtheorem{corollary}{Corollary}[theorem]
\newtheorem*{definition}{Definition}


\newcommand{\norm}[1]{\left\lVert#1\right\rVert}

\newcommand{\Rn}{\mathbb{R}^n}


\newcommand{\too}[1]{\overset{#1}{\longrightarrow}}
\newcommand{\hookar}[1]{\overset{#1}{\hookrightarrow}}
\newcommand{\colim}[1]{ \underset{#1}{\texttt{colim}} \text{ } }
\newcommand{\fset}[1]{ \{ #1_0 , \dots , #1_n \} }
\newcommand{\ddt}{\frac{d}{dt}}


\newcommand{\N}{\mathbb{N}}
\newcommand{\Z}{\mathbb{Z}}






\begin{document}



\setlength{\parindent}{0cm}

\textbf{Ejercicio 1 (12 puntos):} Dada la función $f(x,y)= \frac{\cos(x^3+y^3)}{\ln(x^2+y^2)}$, determinar $f_{xx}$, $f_{yy}$ y $f_{xy}$ \\


\textbf{Solución: }Derivamos

\begin{align*}
    f_x &= \frac{-3x^2\ln(x^2+y^2)\sin(x^3+y^3)-\cos(x^3+y^3)\frac{2x}{x^2+y^2}}{(\ln(x^2+y^2))^2} \\
    f_y &= \frac{-3y^2\ln(x^2+y^2)\sin(x^3+y^3)-\cos(x^3+y^3)\frac{2y}{x^2+y^2}}{(\ln(x^2+y^2))^2}
\end{align*}

Vamos a necesitar definir unas funciones para que no quede tan largo. Sea

\begin{align*}
    A &= -3x^2\ln(x^2+y^2)\sin(x^3+y^3)-\cos(x^3+y^3)\frac{2x}{x^2+y^2} \\
    B &= -3y^2\ln(x^2+y^2)\sin(x^3+y^3)-\cos(x^3+y^3)\frac{2y}{x^2+y^2}
\end{align*}

Entonces

\begin{align*}
    f_{xx} &= \frac{(\ln(x^2+y^2))^2A_x- \frac{4Ax\ln(x^2+y^2)}{x^2+y^2}}{(\ln(x^2+y^2))^4} \\
   f_{yy} &= \frac{(\ln(x^2+y^2))^2B_y- \frac{4By\ln(x^2+y^2)}{x^2+y^2}}{(\ln(x^2+y^2))^4} \\
   f_{xy} &= \frac{(\ln(x^2+y^2))^2A_y- \frac{4Ay\ln(x^2+y^2)}{x^2+y^2}}{(\ln(x^2+y^2))^4} \\
\end{align*}

Hasta aquí se considera bien porque el resulatado es extremadamente largo.


\newpage


\textbf{Ejercicio 2 (12 puntos):} Halle todos los puntos críticos de la función 

\begin{align*}
    f(x,y) &= y^3-x^3-6xy
\end{align*}

Clasifique estos puntos como máximo relativo, mínimo relativo o punto de silla. \\

\textbf{Solución: } Calculamos el gradiente

\begin{align*}
    \nabla f = (-3x^2-6y, 3y^2-6x)
\end{align*}

Luego para tener un punto crítico necesitamos  $(-3x^2-6y, 3y^2-6x)=(0,0)$, luego

\begin{align*}
    -3x^2-6y &=0 \\
    3y^2-6x &=0
\end{align*}

O bien

\begin{align*}
    -x^2-2y &=0 \\
    y^2-2x &=0
\end{align*}

De la segunda ecuación obtenemos que $x = \frac{y^2}{2}$, reemplazando en la primera, obtenemos


\begin{align*}
    -\frac{y^4}{4}-2y &=0 \\
    y^4+8y &=0 \\
    y(y^3+8) &=0 \\
\end{align*}

De donde hay solo dos opciones $y=0$ o bien $y=-2$. Si $y=0$, entonces $x=0$ y si $y=-2$, entonces $x=2$. Por lo tanto los puntos críticos son $(0,0)$ y $(2,-2)$. Para saber que tipo de puntos son calculamos el Hessiano

\begin{align*}
    \begin{vmatrix}
        f_{xx} & f_{xy} \\
        f_{yx} & f_{yy}
    \end{vmatrix} &= \begin{vmatrix}
        -6x & -6 \\
        -6 & 6y
    \end{vmatrix} \\
    &= -36xy-36
\end{align*}

Al evaluar en $(0,0)$ el determinante es menor que cero y por lo tanto es un punto crítico. Por otra parte para $(2,-2)$ se tiene que el determiante es mayor que cero y $f_{xx}(2,-2) =-12$ es menor que cero por lo tanto es un máximo local



\newpage



\textbf{Ejercicio 3 (12 puntos):} Cuando se utilizan $x$ unidades de mano de obra e $y$ unidades de capital, la producción total $P$ de un fabricante está dada por la función de producción:

\begin{align*}
    P(x,y) = 60x^{2/3}y^{1/3}
\end{align*}

Los costos de cada unidad de mano de obra y del capital son de US\$ 64 y \\US\$ 108 por unidad respectivamente. Suponga que el fabricante decide elaborar 2160 unidades de su producto. Halle el número de unidades de mano de obra y de capital que deben de emplearse con objeto de Minimizar el Costo Total. \\

\textbf{Solución: } Queremos minimizar el costo total
\[
f(x,y)=64x+108y
\]
sujeto a la restricción de producción
\[
g(x,y) =60x^{2/3}y^{1/3}-2160.
\]
Según la regla de multiplicadores de Lagrange necesitamos

\[
\nabla f = \lambda \nabla g
\]

Luego 

\begin{align*}
    (64,108) = (40\lambda x^{-1/3}, 20\lambda y^{-2/3})
\end{align*}

De la primera coordenadad obtenemos
\[
40\lambda x^{-1/3}y^{1/3}=64
\quad\Rightarrow\quad
\lambda=-\frac{64}{40}x^{1/3}y^{-1/3}=-\frac{8}{5}x^{1/3}y^{-1/3}.
\]
De la segunda coordenada obtenemos
\[
20\lambda x^{2/3}y^{-2/3}=-108
\quad\Rightarrow\quad
\lambda=-\frac{108}{20}x^{-2/3}y^{2/3}=-\frac{27}{5}x^{-2/3}y^{2/3}.
\]

Igualando las dos expresiones para \(\lambda\)
\[
-\frac{8}{5}x^{1/3}y^{-1/3}=-\frac{27}{5}x^{-2/3}y^{2/3}
\quad\Longrightarrow\quad
8x^{1/3}y^{-1/3}=27x^{-2/3}y^{2/3}.
\]
Multiplicando por \(x^{2/3}y^{1/3}\) en ambos lados obtenemos
\[
8x=27y
\quad\Rightarrow\quad
x=\frac{27}{8}y.
\]

Ahora usamos la restricción \(60x^{2/3}y^{1/3}=2160\), es decir \(x^{2/3}y^{1/3}=36\). Sustituyendo \(x=\tfrac{27}{8}y\):
\[
\left(\frac{27}{8}y\right)^{2/3}y^{1/3}
=\left(\frac{27}{8}\right)^{2/3}y^{2/3}y^{1/3}
=\frac{9}{4}y =36.
\]
De donde
\[
y = 36\cdot\frac{4}{9}=16.
\]
Y entonces
\[
x=\frac{27}{8}\cdot 16 =54.
\]


\newpage



\textbf{Ejercicio 4 (12 puntos):} Considere el campo vectorial $F(x,y) = (e^{2y}, 1+2xe^{2y})$ \\



(a) demuestre que es conservativo \\
(b) calcule su función potencial \\ 
(c) luego calcule la integral de línea sobre $F$: \\

\begin{align*}
    \int_C F \cdot dr 
\end{align*}

Donde $C: r(t) =(te^t,(1+t))$ con $0 \leq t \leq 1$. \\

\textbf{Solución: } a) Tenemos que $M= e^{2y}$ y $N= 1+2xe^{2y}$, tenemos que

\begin{align*}
    N_x=M_y= 2e^{2y}
\end{align*}

POo lo tanto es conservativo. \\


b) Para la función potencia, calculamos

\begin{align*}
    f(x,y) &= \int_0^1 F(tx,ty) \cdot (x,y) dt \\
    &= \int_0^1 (e^{2ty}, 1+2txe^{2ty} )\cdot (x,y) dt \\
    &= \int_0^1 xe^{2ty} + y+2txye^{2ty}  dt \\
    &= \left(\frac{x}{2y}e^{2ty}+ty+xe^{2ty}\left(t-\frac{1}{2y}\right) \right)\Big|_0^1\\
    &= \left(ty+txe^{2ty} \right)\Big|_0^1\\
    &= xe^{2y}+y
\end{align*}

c) Notamos que el punto inicial es $r(0)=(0,1)$ y el fial es $r(1) = (e,2)$, luego por el teorema fundamental de las integrales de línea

\begin{align*}
    \int_C F \cdot dr &= f(e,2) -f(0,1) \\
    &= ee^4+2 -1 \\
    &= e^5+1 \\
\end{align*}


\newpage


\textbf{Ejercicio 5 (12 puntos):} Utilice el Teorema de Green para calcular $\displaystyle \int_C F \cdot dr$, donde $F(x,y) =(y^2-x^2,xy^2)$  y $C$ es el círculo $x^2+y^2=4$ de $(2,0)$ a $(\sqrt{2},\sqrt{2})$  y  los segmentos de recta de $(\sqrt{2},\sqrt{2})$ a $(0,0)$ y de $(0,0)$ a $(2,0)$. \\

\textbf{Solución: } Notemos que $M= y^2-x^2$ y $N= xy^2$, luego $M_y= 2y$ y $N_x= y^2$. Por el teorema de Green tenemos que 

\begin{align*}
    \int_C F \cdot dr = \iint_R y^2-2y dA 
\end{align*}

La región $R$ se puede pasar a polares y nos queda

\begin{align*}
    \int_C F \cdot dr &= \iint_R y^2-2y dA \\
                    &= \int_0^{\pi/4} \int_0^2 (r^2\sin^2(\theta)-2r\sin(\theta))r drd\theta \\
                    &= \int_0^{\pi/4} \int_0^2 r^3\sin^2(\theta)-2r^2\sin(\theta) drd\theta \\
                    &= \int_0^{\pi/4}\left(\sin^2(\theta)\frac{r^4}{4}-\frac{2r^3}{3}\sin(\theta) \right) \Big|_0^2 d\theta \\
                    &= \int_0^{\pi/4} 4\sin^2(\theta)-\frac{16}{3}\sin(\theta)  d\theta \\
                    &= \left(4\left( \frac{\theta}{2}-\frac{\sin(2\theta)}{4}\right)+\frac{16}{3}\cos(\theta)  \right) \Big|_0^{\pi/4}\\
                    &= \frac{\pi}{2}-1 +\frac{8\sqrt{2}}{3}- \frac{16}{3} \\
                    &= \frac{\pi}{2}-\frac{19}{3} +\frac{8\sqrt{2}}{3}
\end{align*}

\end{document}